\chapter{Instalação e execução}

\eyebrow{colocando o instrumento em funcionamento}

\section{Opção A: uma versão pré-compilada}

O caminho mais rápido é baixar um build nativo para a sua plataforma a partir da
página de \textbf{Releases} do projeto. Cada versão distribui um arquivo
autocontido por sistema operacional, com os pesos do modelo já incluídos e sem
necessidade de ferramentas de build:

\begin{center}
\begin{tabularx}{\linewidth}{@{}l X@{}}
\toprule
\textbf{Plataforma} & \textbf{Arquivo} \\
\midrule
macOS (universal) & \tele{carcass-macos.tar.gz} (contém o \tele{.app}) \\
Windows (amd64)   & \tele{carcass-windows.zip} (contém o \tele{.exe}) \\
Linux (amd64)     & \tele{carcass-linux.tar.gz} (contém o binário) \\
\bottomrule
\end{tabularx}
\end{center}

Descompacte-o e execute o aplicativo. Os recursos de análise em Python ainda
precisam de um ambiente Python com as bibliotecas científicas (veja abaixo); todo
o resto funciona de imediato. As versões são produzidas automaticamente por
integração contínua (Estação~14).

\begin{note}[title={Compilando a partir do código-fonte}]
Se você quiser desenvolver, modificar ou compilar para uma plataforma que não esteja
nas versões, siga as Opções~B e~C abaixo.
\end{note}

\section{Pré-requisitos (compilar a partir do código-fonte)}

\begin{center}
\begin{tabularx}{\linewidth}{@{}l l X@{}}
\toprule
\textbf{Componente} & \textbf{Versão} & \textbf{Necessário para} \\
\midrule
Go            & 1.23+ & Compilar e executar o núcleo \\
Node.js       & 18+   & Compilar o front end \\
Wails CLI     & v2    & Desenvolvimento e builds nativos \\
Python        & 3.10+ & Os sidecars de inferência / análise \\
PyTorch, OpenCV, scikit-image, scipy, rembg & --- & Gordura, conformação, monitor ao vivo \\
\bottomrule
\end{tabularx}
\end{center}

Os recursos de núcleo, captura e importação rodam apenas com Go e Node. Os recursos
baseados em modelos (análise de gordura, conformação, monitor ao vivo) exigem
adicionalmente um ambiente Python com as bibliotecas científicas acima.

\section{Opção B: executar a partir do código-fonte}

\begin{lstlisting}
# instala a Wails CLI
go install github.com/wailsapp/wails/v2/cmd/wails@latest

# dependencias python para os sidecars de analise
pip install torch opencv-python scikit-image scipy rembg onnxruntime
\end{lstlisting}

\subsection*{Pesos do modelo}

Os pesos treinados e a tabela de consulta ColorNaming não são distribuídos na
árvore de código-fonte. Coloque-os sob \tele{model/weights/}:

\begin{center}
\begin{tabularx}{\linewidth}{@{}l X@{}}
\toprule
\textbf{Arquivo} & \textbf{Finalidade} \\
\midrule
\tele{fat\_binary.pth}        & Segmentação de gordura (validada, IoU $\sim$0.92) \\
\tele{direct\_class.pth}      & Grau de acabamento (experimental) \\
\tele{eg\_regression.pth}     & Regressão de espessura de gordura (experimental) \\
\tele{joost\_color\_naming.mat} & Tabela de consulta ColorNaming (exigida por todos) \\
\bottomrule
\end{tabularx}
\end{center}

\section{Executando em desenvolvimento}

\begin{lstlisting}
# desenvolvimento com hot-reload (front end + recompilacao do Go a cada mudanca)
wails dev
\end{lstlisting}

Se o \tele{python3} do seu sistema não tiver as bibliotecas científicas, aponte o
aplicativo para um que as tenha:

\begin{lstlisting}
export CARCASS_PYTHON=/path/to/python-with-torch
wails dev
\end{lstlisting}

\begin{note}[title={A câmera e o WebView}]
A captura por webcam e o monitor ao vivo usam a câmera do sistema através do WebView,
portanto precisam rodar dentro do aplicativo (via \tele{wails dev} ou o binário
compilado), não em um navegador comum. A ponte para o núcleo Go só existe dentro
do aplicativo.
\end{note}

\section{Opção C: compilar um aplicativo nativo}

\begin{lstlisting}
wails build -platform darwin/universal   # .app
wails build -platform windows/amd64      # .exe
wails build -platform linux/amd64        # binario
\end{lstlisting}

Para o aplicativo empacotado, os pesos do modelo são copiados para os recursos da
aplicação de modo que os sidecars de análise possam encontrá-los; o script auxiliar
\tele{scripts/build.sh} faz isso automaticamente após um build.
